\pagetitle{One-page reflective justification / course readings}{Why I redesigned a lesson we already taught}
I did not choose \textit{Invisible Invaders} because I needed a convenient example of a science activity. Aastha and I planned and co-facilitated this camp investigation, and I watched the work become more meaningful when young people could connect a question to a place, an observation, and an explanation. The lesson was not empty of equity or student agency. Our original plan already offered ways to speak, point, draw, observe, and return after a pause. I wanted to preserve those commitments while making the science story easier for students to direct and the evidence easier for anyone to challenge.

Our camp record also tells me where our design can improve. We had traceable shoreline samples, a moving-air paper model, and a before/during/after explanation, but those are not three kinds of proof of the same thing. A labeled sand sample is not a confirmed plastic sample. A piece of paper crossing a taped line is evidence about a classroom model, not a measured environmental exposure. Our post-camp reflection recorded requests for more time to investigate and build, less rushing, more movement, and an equivalent sensory-safe route around the fish activity (Garcia \& Aastha, 2026). Those requests matter more to me than whether we finished every planned station.

Reiser et al. (2021) describe the importance of coherence from the learner's perspective. That led me to replace a rapid series of activities with a visible question--investigation--revision arc: students begin with a Friday evidence story, select a comparison, and return to their first model with a reason to change it. Krajcik and Merritt (2012) help explain why the model cannot be a decorative final diagram: its job is to make a tentative account that evidence can refine. Windschitl and Thompson (2013) also connect public modeling and peer questions to making reasoning visible. That is why I added evidence-status labels, two model versions, and space for students to disagree without requiring every uncertainty to disappear.

JuST is not just a final discussion about pollution. Luehmann et al. (2024) push me to consider whose knowledge and questions have authority throughout the investigation, including multilingual, local, and community knowledge. Through my Puzzle Plan and EQUITAS lens, I notice how language, identity, energy, sensory demands, and tool design can interact. I do not want to confuse access friction with scientific capability. A student who draws a pathway, revises a claim in Spanish, records a count, or asks what a sample cannot prove is doing real science. Giving multiple entry routes means little if the teacher still decides every important question.

If I teach this again, I will protect fewer, deeper moves: one student-generated puzzle, one fair comparison, one model revision, and one evidence-based next question or action. I want young people to leave knowing that their ideas changed the explanation---and that saying \emph{we do not know yet} can be a stronger scientific answer than pretending we do.
